\documentclass[11pt]{article}
\usepackage[margin=1in]{geometry}
\usepackage{amsmath,amssymb,booktabs,hyperref}
\newcommand{\dd}{\mathrm d}
\title{Radiative thresholds, vacuum matching, and a gravity compatibility test}
\author{Zimmerman programme: CPB-RV-2026-10-16}
\date{16 October 2026}
\begin{document}
\maketitle
\section{Result and conventions}
The new charge-even mediator has a calculable one-loop tadpole and carrier
mass correction. The small fractional carrier correction can exceed the
original narrow source detuning. A separately specified weaker-coupling
comparison keeps that decay open without finite pole-mass readjustment in
the retained calculation. Neither comparison is a complete renormalized
particle/medium theory or a new observational exclusion.

The vacuum calculation also makes a distinction essential to the full
cold/dark-energy story: fixing a field's vacuum expectation value and mass
does not fix the additive gravitational vacuum energy. A known sequestering
gravity option is examined as a separate route, with its hypotheses and a
nonzero feedback from a proposed acceleration coefficient. No complete new
gravity action, predicted vacuum magnitude or novel mechanism is claimed.

The loop sector uses the same scalar normalization and parameters as the
previous report. The reference renormalization scale is $\mu_R=m$.
The symbol $\Delta_c=M^2-4m^2$ denotes source detuning; $\Delta_f$
continues to mean a nodal gap where needed. They are different quantities.

\section{The scalar one-loop potential and required matching}
On the homogeneous $\Phi=X=0$ axis, a constant mediator background gives
the two real daughter fluctuations the common squared mass
$z(\sigma)=m^2+a\sigma$. For $z>0$, their
$\overline{\rm MS}$ contribution is
\begin{equation}
 V_1^\Phi(\sigma)={z(\sigma)^2\over32\pi^2}
       \left[\ln{z(\sigma)\over\mu_R^2}-{3\over2}\right].
 \label{rv:CW}
\end{equation}
The normalization follows from the standard real-scalar formula in
Ref.~\cite{rv:martin}, Eqs.~(1.1) and (3.2), with two real degrees of
freedom. This is the scalar-only contribution on this axis, not the full
Higgs, medium, gauge or two-loop effective potential. A region with
$z\leq0$ is outside this real-axis expansion and cannot be declared stable
by extrapolating its Taylor series.

The derivatives at the origin are
\begin{align}
 T_\sigma&={a m^2\over16\pi^2}
                         \left[\ln{m^2\over\mu_R^2}-1\right],\\
 \delta M_{\sigma,0}^2&={a^2\over16\pi^2}\ln{m^2\over\mu_R^2},\\
 V_1'''(0)&={a^3\over16\pi^2m^2},\qquad
 V_1''''(0)=-{a^4\over16\pi^2m^4}.
\end{align}
The last two are finite zero-momentum vertices in this calculation, not
evidence that independent cubic/quartic UV counterterms are required by
this diagram. The negative fourth Taylor coefficient alone does not prove
an instability of the exact multiscalar potential.

For an independent UV-operator check write
$\Phi=(f_1+if_2)/\sqrt2$, $X=(x_1+ix_2)/\sqrt2$ and form the real
five-field potential Hessian $\mathcal H$. With
$\lambda_\Phi=\lambda_{\rm tot}+\delta\lambda$, its
$\sigma$-dependent squared trace is exactly
\begin{equation}
 \operatorname{Tr}\mathcal H^2(\sigma)
 -\operatorname{Tr}\mathcal H^2(0)
 =2a^2\sigma^2+4am^2\sigma+16a\lambda_\Phi\sigma|\Phi|^2.
\end{equation}
Thus a tadpole counterterm is needed even though the tree construction
chose $\sigma=0$ at empty vacuum. The same normalization gives the retained
one-loop contributions
\begin{equation}
 \beta_{t_\sigma}={am^2\over8\pi^2},\quad
 \beta_{M_\sigma^2}={a^2\over8\pi^2},\quad
 \beta_a={a\lambda_\Phi\over2\pi^2},\quad
 \left.\beta_{m^2}\right|_a={a^2\over8\pi^2}.
 \label{rv:betas}
\end{equation}
These are selected scalar contributions, not the complete beta-function
system. At this order the explicit scale derivative of
Eq.~\eqref{rv:CW} is cancelled on the axis by
\begin{equation}
 \beta_{V_0}^\Phi={m^4\over16\pi^2},\qquad
 \beta_{t_\sigma}\sigma+{\beta_{M_\sigma^2}\over2}\sigma^2.
\end{equation}

At $\mu_R=m$, keeping the renormalized tree linear coefficient $t_R=0$
without an additional tadpole-matching condition shifts the minimum at leading order by
$\delta\sigma=-T_\sigma/M_\sigma^2$.
For $a=.0004$ GeV and $M_\sigma=.02$ GeV this is
$2.53303\,10^{-8}$ GeV. Its additional fractional carrier mass-squared
shift is $a\delta\sigma/m^2=2.53303\,10^{-6}$. An alternative, used in
the next pole comparison, fixes $\langle\sigma\rangle=0$ by a finite
tadpole condition. These are different matching choices; neither fixes
the vacuum constant or all other masses.

\section{The narrow source threshold is the stronger test}
With the tadpole separately fixed to zero, the new mediator's one-loop
contribution to the daughter pole mass at $p^2=m^2$ is
\begin{equation}
 \delta m_{\Phi,\rm pole}^2
 ={a^2\over16\pi^2}\int_0^1\dd x\,
       \ln{xM_\sigma^2+(1-x)^2m^2\over\mu_R^2}.
 \label{rv:pole}
\end{equation}
This is the finite $\overline{\rm MS}$-parameter comparison at the stated
scale, not a prediction obtained while simultaneously fixing the pole mass
to its old numerical value. The new $a$ interaction does not shift the
$X$ pole at this retained one-loop order. Consequently
\begin{equation}
 \delta\Delta_{c,\rm pole}=-4\delta m_{\Phi,\rm pole}^2,
 \qquad
 \left.\beta_{\Delta_c}\right|_a=-{a^2\over2\pi^2}.
\end{equation}
For the original $m=2$ MeV, $\Delta_c=6.4\,10^{-9}$ GeV$^2$ family:
\begin{center}
\begin{tabular}{rrrrr}
\toprule
$M_\sigma$ (MeV)&$a$ (MeV)&$\delta m_\Phi^2/m^2$
 &$\Delta_{c,\rm pole}$ (GeV$^2$)&decay open\\
\midrule
6 & .12 & $3.2631\,10^{-5}$ & $5.87790\,10^{-9}$ & yes\\
10 & .20 & $1.4793\,10^{-4}$ & $4.03307\,10^{-9}$ & yes\\
20 & .40 & $9.2379\,10^{-4}$ & $-8.38071\,10^{-9}$ & no\\
200 & 4.00 & $.20799$ & $-3.32148\,10^{-6}$ & no$^*$\\
20 & .10 & $5.7737\,10^{-5}$ & $5.47621\,10^{-9}$ & yes\\
\bottomrule
\end{tabular}
\end{center}
All entries retain only the new $a^2$ correction with the old numerical
masses treated as renormalized inputs at $\mu_R=m$. The starred case also
fails the declared $|\delta m^2|/m^2<0.1$ loop-size diagnostic; its
quantitative extrapolation is therefore not certified. A successful
numerical integration does not override that scope failure.

The threshold diagnostic is scale-consistent. Running the input
$\Delta_c$ with Eq.~\eqref{rv:betas} cancels the explicit logarithm in
Eq.~\eqref{rv:pole} at the retained order. Changing the scale while
holding all those renormalized numerical inputs fixed would instead define
a different physical matching choice. Alternatively one may impose the
original chosen pole masses as inputs and choose the finite counterterms
accordingly. The negative-gap entry does not exclude that on-shell model;
it demonstrates the threshold matching it needs.
The retained cancellation is not a complete RG flow or a certificate that
the tree positive-square parameterization stays in its positive cone at
every scale.

At $M_\sigma=20$ MeV the no-readjustment criterion
$4\delta m_\Phi^2<\Delta_c$ gives $a<.263210$ MeV in this comparison.
The separately specified $a=.10$ MeV continuation lies below it. Its
compensated tree coefficients are
\begin{equation}
 \delta\lambda=1.25\,10^{-5},\quad
 \zeta_{\rm NR}=.625\,\mathrm{GeV}^{-2},\quad
 u_R=.000376103\,\mathrm{GeV}^{-2},\quad
 u_0=.250376\,\mathrm{GeV}^{-2}.
\end{equation}
It retains the symbolic tree positive-square condition and passes the
stated new-loop mass/threshold comparisons. Its thermodynamics, actual
abundance and full renormalized amplitudes are not inherited from the
older stronger candidate.

The mediator's own vacuum bubble supplies an independent normalization
check:
\begin{equation}
 \delta\mathcal K_\sigma(p^2)
 ={a^2\over16\pi^2}\int_0^1\dd x\,
 \ln{m^2-x(1-x)(p^2+i0)\over\mu_R^2}.
\end{equation}
At $p^2=M_\sigma^2>4m^2$, its imaginary part obeys
$-\operatorname{Im}\delta\mathcal K_\sigma/M_\sigma
=\Gamma_{\sigma\to\Phi\Phi^*}$, reproducing the preceding tree width.
At low $p^2$ its induced kinetic coefficient is
$\delta Z_\sigma=a^2/(96\pi^2m^2)>0$.
A zero curvature correction at $\mu_R=m$ therefore does not imply a zero
on-shell mediator self-energy.
This vacuum calculation is different from absorption into the fermion
medium, whose density vertex has the uniform Ward zero found earlier.

\section{Vacuum energy is an independent matching condition}
At the chosen reference scale the free scalar fluctuation pieces are
\begin{align}
 V_1^\Phi(0)&=-{3m^4\over64\pi^2},\\
 V_1^X(0)&={M^4\over32\pi^2}
                   \left[\ln{M^2\over m^2}-{3\over2}\right],\\
 V_1^\sigma(0)&={M_\sigma^4\over64\pi^2}
                   \left[\ln{M_\sigma^2\over m^2}-{3\over2}\right].
\end{align}
For the 20-MeV mediator these are respectively
$-7.59909\,10^{-14}$, $-9.19156\,10^{-14}$ and
$7.86549\,10^{-10}$ GeV$^4$. A declared comparison reference
$\epsilon_{\rm ref}=(2.25\,\mathrm{meV})^4=2.56289\,10^{-47}$ GeV$^4$
is smaller than their sum by a factor about $3.07\,10^{37}$.
This reference is an input scale, not a new cosmological fit.

Those partial, scheme-dependent terms are not a prediction that the
Universe has that large or signed cosmological constant. The full theory
contains a renormalized vacuum parameter:
\begin{equation}
 \epsilon_{\rm vac,phys}=V_{0,R}(\mu_R)
               +V_{\rm loops}(\hbox{renormalized vacuum};\mu_R).
\end{equation}
Its scale dependence cancels when all retained terms are matched
consistently. Fixing the tadpole, fixing pole masses and fixing this value
are independent conditions. The massive mediator and the nodal medium
have not provided a symmetry or constraint fixing the last one.
Other matter, counterterms, the occupied-state stress and gravity are
absent from this partial scalar sum. No cancellation by the filled bands
can be assumed from a chosen zero of their homogeneous energy.

\section{A distinct known gravity route and its compatibility gate}
Vacuum-energy sequestering is a known possible gravity modification,
not a mechanism invented here. In the matter-loop, semiclassical-gravity
setting of Ref.~\cite{rv:sequester}, nonlinear four-form couplings enforce
constant gravitational auxiliaries. Its Eq.~(10), in that source's gravity
conventions, is
\begin{equation}
 M_G^2G^\mu{}_{\nu}
 =\mathcal P[T]^\mu{}_{\nu}-\Delta\Lambda\,\delta^\mu{}_{\nu},
 \qquad
 \mathcal P[T]^\mu{}_{\nu}
 =T^\mu{}_{\nu}-{1\over4}\langle T^\alpha{}_{\alpha}\rangle
                                      \delta^\mu{}_{\nu}.
\end{equation}
The normalized historic spacetime average makes $\mathcal P$ annihilate
a constant multiple of the identity. Residual flux/history data are still
needed, and graviton loops and ill-defined infinite-volume averages are
outside the imported statement. The source requires a covariant matter
action; our preferred-frame band/frame sector has not supplied one.
This remains a conditional route, not an adopted completion of the model.

Three compatibility conclusions can already be checked. First, cancelling
constant vacuum stress does not cancel the local $\sigma$ tadpole or
threshold derivatives. Second, the positive physical residual dark-energy
coefficient $B_{\rm eff}$ need not equal the auxiliary bare $\Lambda$:
it also contains the nonvacuum historic trace contribution after the
vacuum split. Third, a target $a_0^2=\kappa_a^2G B_{\rm eff}$ in natural
units cannot be replaced by a function of the shifting bare counterterm
without deriving the consequences. The source's $M_G$ is not the fitted
acceleration coefficient $\kappa_a$.

For an explicit check, suppose instead that a positive auxiliary vacuum
parameter $B$ is put directly into a toy deep-force energy,
\begin{equation}
 a_0(B)=\kappa_a\sqrt{GB},\qquad
 W(g;B)={g^3\over12\pi G a_0(B)}.
\end{equation}
This is a chosen local deep-force functional, not the solved full nodal
action. Its unavoidable derivative is
\begin{equation}
 \partial_B W=-{W\over2B},\qquad
 {W\over B}={\kappa_a^2\over12\pi}(g/a_0)^3.
\end{equation}
The usual four-form variation of a volume term $-B$ now receives
$-\partial_BW$ as well: its volume integrand is proportional to
$1+\partial_BW$, not simply one. Thus the original constraint and its
vacuum-cancellation argument cannot be carried over unchanged after this
force coupling. For $\kappa_a=1/2$ the correction magnitude is
$(g/a_0)^3/(96\pi)$; it is small in the deep regime but not zero.
This does not disprove a completed residual-flux coupling. It identifies
the extra equation such a construction must solve. A bare-counterterm
identification, or an unchanged sequestering equation with a new
$B$-dependent matter action inserted afterwards, is insufficient.

\section{What is established and what remains}
The scalar bridge has concrete radiative matching obligations and a
weaker-coupling continuation which survives the retained threshold test.
The reference scalar vacuum calculation exposes the independent vacuum
condition rather than predicting it. A known gravity route can address
constant matter-loop stress under additional hypotheses; the acceleration
coupling changes those hypotheses unless its auxiliary variation is counted.
This is progress toward a joined action, not whole-theory closure.

Next are the complete matching of tadpoles, masses, interactions and the
open positive region, the neutral mixture EOS and populated-mediator
kinetics, and an explicit covariant frame/residual-vacuum action if the
gravity route is pursued. Finite supply, physical boundaries, formation,
source selection, metric/lensing, native cosmology and joint empirical
gates remain open. No unique or globally novel mechanism is established.

Evidence: \texttt{scalar\_radiative\_audit\_run\_003} has 33 controls,
including the retained RG gap check and a weaker comparison;
\texttt{vacuum\_bridge\_compatibility\_run\_001} has seven algebraic
controls. The first radiative run failed its original quadrature gate for
$M_\sigma/m=100$. Refinement from 128 to 512 Gauss points fixes the
endpoint resolution without relaxing the $10^{-10}$ tolerance. Run 002
completed the calculation; run 003 adds the scale-consistency checks.
Physical scope failures remain in the output. Review is self-review.

\begin{thebibliography}{9}
\bibitem{rv:martin} S. P. Martin, Two-loop effective potential for a general
renormalizable theory and softly broken supersymmetry,
\href{https://arxiv.org/pdf/hep-ph/0111209v2}{hep-ph/0111209v2},
Phys. Rev. D 65 (2002) 116003, Eqs.~(1.1), (3.2).
\bibitem{rv:sequester} N. Kaloper, A. Padilla, D. Stefanyszyn and
G. Zahariade, A Manifestly Local Theory of Vacuum Energy Sequestering,
\href{https://arxiv.org/pdf/1505.01492v2}{arXiv:1505.01492v2},
Phys. Rev. Lett. 116 (2016) 051302, Eqs.~(7)--(10).
\end{thebibliography}
\end{document}
